\subsubsection{arity-charge 的大偏差率函数（\texorpdfstring{$\chi$}{chi} 切片）}\label{subsec:chi-ldp}
沿 $\chi$-切片 $\theta(t)=(t,0,-t)$，压力函数 $P_\chi(t)=\log\rho(M_{\theta(t)})$ 给出 $\chi$ 的指数型倾斜族；相应的大偏差率函数可由 Legendre 对偶直接读出。

\begin{remark}[标准定理接口：Gärtner--Ellis]\label{thm:chi-gartner-ellis}
对原始矩阵族的 \(\chi\)-切片压力 \(P_\chi(t)=\log\rho(M_{\theta(t)})\)，其 Legendre 对偶
\(
I_\chi(\alpha)=\sup_{t\in\RR}(t\alpha-P_\chi(t))
\)
给出经验均值的大偏差率函数；并在 \(\bar\alpha=P_\chi'(0)\) 处有二次近似 \(I_\chi(\bar\alpha+\delta)=\delta^2/(2P_\chi''(0))+o(\delta^2)\)。参见 \cite{DemboZeitouni2010LargeDeviations}。
\end{remark}

\begin{remark}[\texorpdfstring{$\chi$}{chi} 的支撑端点与端点压力（真实输入 40 状态核）]\label{rem:chi-ldp-endpoints}
对真实输入 $40$ 状态核的 essential $20$ 状态核心，附录 \ref{app:real-input-40-arity-charge} 的定理
\ref{thm:real-input-40-arity-charge-density-bound}
给出对任意 primitive 周期轨道 $\gamma$（$n=|\gamma|$）的硬约束
\(
0\le \chi(\gamma)\le \lfloor n/2\rfloor
\)，从而任意经验平均
\(
\bar\chi_n=\chi(\gamma)/n
\)
必落在区间 $[0,1/2]$。
进一步，附录同节命题 \ref{prop:real-input-40-arity-charge-endpoints} 给出端点压力：
当 $t\to-\infty$（即 $q=e^t\to 0^+$）时，
\[
P_\chi(t)\to \log\kappa,
\qquad
\kappa^4-\kappa^3-1=0,\ \kappa>1;
\]
当 $t\to+\infty$（即 $q\to+\infty$）时，
\[
P_\chi(t)=\frac12\,t+\frac12\log 3+o(1).
\]
因此 Legendre 对偶率函数 $I_\chi(\alpha)=\sup_t(t\alpha-P_\chi(t))$ 在 $\alpha\notin[0,1/2]$ 时必取 $+\infty$（由端点线性斜率直接推出），并且其有效支撑端点正由
\(
\lim_{t\to-\infty}P_\chi'(t)=0
\)
与
\(
\lim_{t\to+\infty}P_\chi'(t)=1/2
\)
钉住。
\end{remark}

\begin{proposition}[\texorpdfstring{$\chi$}{chi}-扭曲 \texorpdfstring{$\zeta$}{zeta} 的 $7$ 次谱接口（真实输入 40 状态核）]\label{prop:chi-degree7-interface}
\leanverified{paper\_chi\_degree7\_interface}
对真实输入 \(40\) 状态核的 essential \(20\) 状态核心，沿 \(\chi\)-切片令 \(q:=e^t>0\)，并记 \(M(q):=M_{\theta(t)}\)（等价地：对每条边赋权 \(q^{\chi}\) 得到的加权矩阵）。
则
\[
\zeta_\chi(z,q):=\frac{1}{\det(I-zM(q))}
\]
为有理函数，并且其分母具有显式因子分解（脚本生成）：
% AUTO-GENERATED by scripts/exp_sync_kernel_real_input_40_arity_charge_closed_form.py
\[
\begin{aligned}
\det(I-zM(q))=(1+z)(1-qz^2)\,Q_7(z,q),
\\
Q_7(z,q)&=1 - 2\,z + \left(1 - 5 q\right)\,z^{2} + 6 q\,z^{3}\\
&+ \left(6 q^{2} - 3 q - 1\right)\,z^{4} + \left(- 4 q^{2} - q + 1\right)\,z^{5}\\
&+ q \left(4 q - 3\right)\,z^{6} + q \left(1 - q\right)\,z^{7}.
\\[2pt]
\det(\Lambda I-M(q))=\Lambda^{10}(\Lambda+1)(\Lambda^2-q)\,P_7(\Lambda,q),
\\
P_7(\Lambda,q)&=q \left(1 - q\right) + q \left(4 q - 3\right)\,\Lambda\\
&+ \left(- 4 q^{2} - q + 1\right)\,\Lambda^{2}\\
&+ \left(6 q^{2} - 3 q - 1\right)\,\Lambda^{3} + 6 q\,\Lambda^{4}\\
&+ \left(1 - 5 q\right)\,\Lambda^{5} - 2\,\Lambda^{6} + \Lambda^{7}.
\end{aligned}
\]

特别地，除去平凡谱支 \(\Lambda=-1\) 与 \(\Lambda=\pm\sqrt q\) 外，全部非零非平凡本征值由次数 \(7\) 的代数曲线 \(P_7(\Lambda,q)=0\) 控制；并且分母只含 \(q\) 的非负幂（\(Q_7\in\ZZ[z,q]\)）。
\end{proposition}

\begin{remark}[与 Abel 有限部常数曲线的同一切片对比]
$P_\chi(t)$ 控制 leading 指数尺度（典型密度、涨落与大偏差），而 $\log\mathfrak{M}_\chi(t)$ 控制 prime-orbit Mertens 主项中的常数级漂移：
\[
\sum_{n\le N}\frac{p_{n,\chi}(t)}{\lambda_\chi(t)^n}=\log N+\gamma+\log\mathfrak{M}_\chi(t)+o(1),
\qquad
\lambda_\chi(t)=e^{P_\chi(t)}.
\]
因此在同一切片上同时记录 $t\mapsto(P_\chi(t),\log\mathfrak{M}_\chi(t))$，可把“指数尺度的代价”（由 $P_\chi$ 给出）与“常数级漂移的再分配”（由 $\log\mathfrak{M}_\chi$ 给出）作为两种互补指纹并列比较。
在真实输入核上，命题 \ref{prop:chi-degree7-interface} 进一步将该切片的谱/极点接口压缩为一个显式的 \(7\) 次代数曲线：因此 \(\lambda_\chi(t)\)、主极点留数常数 \(C_\chi(t)\) 以及 \(\log\mathfrak{M}_\chi(t)\) 的 Möbius--Abel 尾项，均可在该曲线及其导数上完成计算，而无需回到 \(40\times 40\) 矩阵层。
并且由端点渐近 \(P_\chi(t)=\tfrac12t+\tfrac12\log 3+o(1)\)（注 \ref{rem:chi-ldp-endpoints}）知 \(z_\star(t)=e^{-P_\chi(t)}\asymp e^{-t/2}\)，从而对任意固定 \(k\ge 2\)，采样点 \(z_\star(t)^k\) 随 \(t\to+\infty\) 以 \(e^{-kt/2}\) 速度趋零，Möbius 尾项收敛极快。
\end{remark}

\paragraph{内生时间尺度（无外部输入）：\texorpdfstring{$\chi$}{chi} 切片的相关时间}
本小节的 $t$ 参数不仅控制指数倾斜的静态代价（$P_\chi(t)$）与常数级漂移（$\log\mathfrak{M}_\chi(t)$），也同时诱导一条\textbf{纯核内生}的时间尺度：在 $\theta(t)=(t,0,-t)$ 的平衡态下，arity-charge 观测量的时间相关以谱隙速率指数衰减。
为避免与角色投影 $P_\chi$ 的混合时间 $\tau_{\mathrm{mix}}$（注 \ref{rem:pom-chi-mixing}）混淆，这里把该尺度记为 $\tau_{\mathrm{corr}}$。

\begin{remark}[\texorpdfstring{$\chi$}{chi} 切片的相关时间（谱隙口径）]\label{prop:chi-slice-time-correlation}
设真实输入 40 状态核在 $\theta(t)=(t,0,-t)$ 下的加权矩阵为 $M_{\theta(t)}$，并假设在给定 $t$ 的邻域内矩阵族原始。
定义
\[
\lambda_1(t):=\rho\bigl(M_{\theta(t)}\bigr),
\qquad
\Lambda(t):=\max\{ |\lambda|:\ \lambda\in\mathrm{spec}(M_{\theta(t)}),\ |\lambda|<\lambda_1(t)\},
\qquad
\rho_{\mathrm{corr}}(t):=\frac{\Lambda(t)}{\lambda_1(t)}\in[0,1),
\]
以及相关时间
\[
\tau_{\mathrm{corr}}(t):=\frac{1}{-\log \rho_{\mathrm{corr}}(t)}.
\]
对有限状态原始马尔可夫链，中心化有限窗口可观测量的时间相关满足指数衰减包络
\(
|\mathrm{Cov}(\chi_0,\chi_n)|\le C(t)\,\rho_{\mathrm{corr}}(t)^n
\)，其中 $C(t)<\infty$（标准谱隙结论；参见 \cite{LevinPeresWilmer2009MarkovMixing}）。
本文仅使用其“由有限维谱读出指数包络”的可审计接口。
\end{remark}

\input{sections/generated/fig_real_input_40_time_correlation}

\begin{remark}[黄金点的闭式相关时间（真实输入 40 状态核）]\label{cor:chi-slice-t0-golden-tau}
对真实输入 40 状态核在 $t=0$ 的 $\chi$ 切片平衡态，谱隙比满足
\[
\rho_{\mathrm{corr}}(0)=\frac{\Lambda(0)}{\lambda_1(0)}=\varphi^{-1},
\qquad
\tau_{\mathrm{corr}}(0)=\frac{1}{-\log\rho_{\mathrm{corr}}(0)}=\frac{1}{\log\varphi}.
\]
该闭式与图 \ref{fig:real-input-40-time-correlation} 的 $t=0$ 读出一致。
\end{remark}

\paragraph{事件时钟与时间重标定（无外部输入）}
在 $t=0$ 时，平衡态退化为未加权核的 Parry 测度，因此同一条平衡态轨道上可以引入一个\textbf{独立于 $\chi$ 的“事件时钟”}：令 $e_t\in\{0,1\}$ 表示真实输入核在第 $t$ 步输出位（见附录中对输出位势的定义口径），并定义事件时间
\[
T_n:=\sum_{t=0}^{n-1} e_t.
\]
则由输出位势的压力函数 $P_e(\theta)=\log\lambda(e^{\theta})$，其均值密度为 $\mu_e:=\mathbb{E}[e_0]=P_e'(0)$（闭式可审计，见表 \ref{tab:real_input_40_output_potential_cumulants_closed}）。
将命题 \ref{prop:chi-slice-time-correlation} 的“步数时间”相关尺度 $\tau_{\mathrm{corr}}(0)$ 通过 $T_n\approx \mu_e n$ 换元，即得到以事件时间计量的等价尺度 $\tau_{\mathrm{corr}}^{(e)}:=\mu_e\tau_{\mathrm{corr}}(0)$。
表 \ref{tab:real-input-40-event-time-dilation} 给出该换元在 $t=0$ 的数值化读出，作为“内生时间单位选择”对尺度的可审计重标定。

% AUTO-GENERATED by scripts/exp_real_input_40_event_time_dilation.py
\begin{table}[H]
\centering
\small
\setlength{\tabcolsep}{8pt}
\renewcommand{\arraystretch}{1.15}
\caption{Event-time dilation at the unweighted (t=0) Parry equilibrium for the real-input 40-state kernel. We convert the step-time correlation scale $\tau_{\mathrm{corr}}$ to event-time units using the output-bit rate $\mu_e=P_e'(0)$ from the closed-form output potential curve.}
\label{tab:real-input-40-event-time-dilation}
\begin{tabular}{l r}
\toprule
quantity & value\\
\midrule
$r_{\mathrm{corr}}(0)$ & 0.618033988750\\
$-\log r_{\mathrm{corr}}(0)$ & 0.481211825060\\
$\tau_{\mathrm{corr}}$ (step) & 2.078086921235\\
$t_{1/2}$ (step) & 1.440420090413\\
\midrule
$\mu_e=P_e'(0)$ & 0.367376207875\\
$\tau_{\mathrm{corr}}$ (event) & 0.763439692758\\
$t_{1/2}$ (event) & 0.529176070563\\
\bottomrule
\end{tabular}
\end{table}


\begin{remark}[事件时钟下的尺度分离放大（恒等式）]\label{cor:chi-event-clock-amplifies-separation}
在 $t=0$ 的 Parry 平衡态下，令事件时钟 $T_n=\sum_{t=0}^{n-1}e_t$，并令 $\mu_e:=\mathbb{E}[e_0]=P_e'(0)$。
将步数时间相关尺度 $\tau_{\mathrm{corr}}(0)$ 重标定为事件时间尺度
\[
\tau_{\mathrm{corr}}^{(e)}:=\mu_e\,\tau_{\mathrm{corr}}(0).
\]
则角色扭曲混合尺度 $\tau_{\mathrm{mix}}$ 与事件时间相关尺度的比值满足恒等式
\[
\frac{\tau_{\mathrm{mix}}}{\tau_{\mathrm{corr}}^{(e)}}
=\frac{\tau_{\mathrm{mix}}}{\tau_{\mathrm{corr}}(0)}\cdot\frac{1}{\mu_e}.
\]
对真实输入 40 状态核，表 \ref{tab:real-input-40-event-clock-vs-tau-mix} 给出在 $((3,3,5))$ 最坏扭曲处该比值的数值读出，显示事件时钟重标定并不能消除角色扭曲的混合瓶颈，反而将两者的分离进一步放大为一个数量级以上的可审计差异。
\end{remark}

% AUTO-GENERATED by scripts/exp_real_input_40_event_clock_vs_tau_mix.py
\begin{table}[H]
\centering
\small
\setlength{\tabcolsep}{8pt}
\renewcommand{\arraystretch}{1.15}
\caption{事件时钟尺度与角色扭曲混合尺度的分离（真实输入 40 状态核；$t=0$ Parry 平衡态）。事件时间尺度 $\tau_{\mathrm{corr}}^{(e)}=\mu_e\tau_{\mathrm{corr}}$ 来自输出位事件率 $\mu_e=P_e'(0)$ 对步数时间的重标定；而 $\tau_{\mathrm{mix}}$ 来自 $((3,3,5))$ 最坏角色扭曲的谱半径比 $r=\rho_{3,3,5}/\lambda$。}
\label{tab:real-input-40-event-clock-vs-tau-mix}
\begin{tabular}{l r}
\toprule
quantity & value\\
\midrule
$\mu_e=P_e'(0)$ & 0.367376207875\\
$\tau_{\mathrm{corr}}$ (step) & 2.078086921235\\
$\tau_{\mathrm{corr}}^{(e)}$ (event) & 0.763439692758\\
\midrule
$r_{\mathrm{mix}}=\rho_{3,3,5}/\lambda$ & 0.950621088873\\
$\tau_{\mathrm{mix}}$ & 19.747340579192\\
$t_{1/2,\mathrm{mix}}$ & 13.687813446024\\
\midrule
$\tau_{\mathrm{mix}}/\tau_{\mathrm{corr}}^{(e)}$ & 25.866274397982\\
$t_{1/2,\mathrm{mix}}/t_{1/2,\mathrm{corr}}^{(e)}$ & 25.866274397982\\
\bottomrule
\end{tabular}
\end{table}


\begin{remark}[事件时钟的有限样本偏差证书（二阶矩接口）]\label{prop:chi-event-clock-deviation-certificate}
在 $t=0$ 的 Parry 平衡态下，令 $e_t\in\{0,1\}$ 为输出位，$T_n:=\sum_{t=0}^{n-1}e_t$，并记 $\mu_e:=\mathbb{E}[e_0]$。
将一步输出位视为有限状态 Parry 马尔可夫链上的有界可观测量，则协方差序列
\(
\mathrm{Cov}(e_0,e_k)
\)
与部分和方差
\(
\mathrm{Var}(T_n)
\)
均可由有限维转移矩阵 $P$ 的幂与内积\emph{精确}计算。
进一步地，令
$$
S_n:=\sum_{t=1}^{n}\mathrm{Var}(T_t),\qquad b_n(\delta):=\sqrt{\frac{S_n}{\delta}}.
$$
则由并合界与 Chebyshev 不等式得到可审计的统一偏差证书：对任意 $\delta\in(0,1)$ 与任意 $n\ge 1$，有
$$
\mathbb{P}\Bigl(\max_{1\le t\le n}|T_t-\mu_e t|\ge b_n(\delta)\Bigr)\le \delta.
$$
表 \ref{tab:real-input-40-event-clock-deviation-certificate} 给出对真实输入 40 状态核的 $(\mu_e,\mathrm{Var}(T_n),S_n,b_n)$ 的数值化读出；由于其输入仅为核的有限状态线性代数对象，故该证书在本文口径下可完全再现。
\end{remark}

\begin{remark}[更强指数型证书的外部定理接口]
命题 \ref{prop:chi-event-clock-deviation-certificate} 采用二阶矩与并合给出统一界，其优点是无需任何额外假设（例如可逆性）且全部量可由有限维线性代数闭式审计。
若进一步希望得到与混合强度相关的\emph{指数型}尾界，可将 $T_n$ 视为有限状态马尔可夫链的加性泛函，并调用马尔可夫链浓缩不等式（例如 \cite{Paulin2015MarkovConcentration} 的 McDiarmid 型界；在可逆链情形亦可用最优 Hoeffding 型界 \cite{LeonPerron2004OptimalHoeffding}）。
这些外部定理可把 $|T_n-\mu_e n|$ 与对应击中时间 $\tau_k:=\inf\{n:\,T_n\ge k\}$ 的偏差统一写为谱隙/混合参数控制的显式指数尾，从而把“事件时钟重标定”的误差预算进一步收紧到可与 $\tau_{\mathrm{mix}}$ 同尺度比较的形式。
\end{remark}

% AUTO-GENERATED by scripts/exp_real_input_40_event_clock_deviation_certificate.py
\begin{table}[H]
\centering
\scriptsize
\setlength{\tabcolsep}{6pt}
\renewcommand{\arraystretch}{1.15}
\caption{事件时钟偏差的可审计证书（真实输入 40 状态核；$t=0$ Parry 平衡态）。设 $T_n=\sum_{t=0}^{n-1} e_t$ 且 $\mu_e=0.367376207875$。表中 $V_n=\mathrm{Var}(T_n)$ 由有限维 Parry 链精确计算；$S_n=\sum_{t=1}^n V_t$，并定义 $b_n(\delta)=\sqrt{S_n/\delta}$（此处 $\delta=0.05$）。由并合界 + Chebyshev 不等式可得 $\mathbb{P}(\max_{1\le t\le n}|T_t-\mu_e t|\ge b_n(\delta))\le\delta$。此外给出渐近方差密度 $P_e''(0)$ 作为对照。}
\label{tab:real-input-40-event-clock-deviation-certificate}
\begin{tabular}{r r r r r r r}
\toprule
$n$ & $\mu_e n$ & $V_n$ & $V_n/n$ & $P_e''(0)$ & $|V_n/n-P_e''(0)|$ & $b_n(\delta)$\\
\midrule
3 & 1.102128623625 & 0.366849322106 & 0.122283107369 & 0.048741245618 & 0.073541861750 & 4.262015623314\\
10 & 3.673762078751 & 0.702544030613 & 0.070254403061 & 0.048741245618 & 0.021513157443 & 9.795660911793\\
20 & 7.347524157501 & 1.189973354826 & 0.059498667741 & 0.048741245618 & 0.010757422123 & 17.031807898127\\
50 & 18.368810393754 & 2.652210724504 & 0.053044214490 & 0.048741245618 & 0.004302968872 & 38.175385762274\\
100 & 36.737620787507 & 5.089273005421 & 0.050892730054 & 0.048741245618 & 0.002151484436 & 73.160594351536\\
200 & 73.475241575015 & 9.963397567258 & 0.049816987836 & 0.048741245618 & 0.001075742218 & 143.017077246624\\
\bottomrule
\end{tabular}
\end{table}


\begin{remark}[事件时间的粗换时：指数衰减率的 $\mu_e$ 重标定]\label{lem:chi-event-time-rescaling}
在命题 \ref{prop:chi-slice-time-correlation} 的设定下，设某中心化有限窗口可观测量的步数时间相关满足指数包络
$$
\bigl|\mathrm{Cov}(F,\,G\circ\sigma^n)\bigr|\le C\,\rho^n\qquad (n\ge 0),
$$
其中 $\rho\in(0,1)$。
在 $t=0$ 的 Parry 平衡态下引入事件时钟 $T_n=\sum_{s=0}^{n-1}e_s$ 并记 $\mu_e=\mathbb{E}[e_0]>0$。
定义一个仅依赖均值密度的“事件时间近似换时”
$$
\widetilde\tau(k):=\left\lfloor\frac{k}{\mu_e}\right\rfloor\qquad (k\in\mathbb{N}).
$$
则对应的事件时间相关包络满足
$$
\bigl|\mathrm{Cov}(F,\,G\circ\sigma^{\widetilde\tau(k)})\bigr|
\le C\,\rho^{\lfloor k/\mu_e\rfloor}
\le C\,\exp\!\left(-\frac{-\log\rho}{\mu_e}\,k\right).
$$
该结论只使用“步数时间的指数包络”与事件密度 $\mu_e$，从而把步数时间尺度 $\tau_{\mathrm{corr}}=1/(-\log\rho)$ 纯粹地重标定为事件时间尺度 $\tau_{\mathrm{corr}}^{(e)}=\mu_e\tau_{\mathrm{corr}}$。
作为背景：对有限状态可逆马尔可夫链，上述指数包络可由谱隙形式的协方差衰减直接推出（参见 \cite[Chapter~12]{LevinPeresWilmer2009MarkovMixing}）。
\end{remark}

\begin{remark}[击中时间 $\tau(k)$ 的可审计偏差证书（方差型）]
\label{prop:chi-event-clock-hitting-time-certificate}
\leanverified{paper\_chi\_event\_clock\_hitting\_time\_certificate}
在 $t=0$ 的 Parry 平衡态下，令 $e_t\in\{0,1\}$ 为输出位、$T_n:=\sum_{t=0}^{n-1}e_t$、$\mu_e:=\mathbb{E}[e_0]>0$，并定义击中时间
$$
\tau(k):=\inf\{n\ge 1:\,T_n\ge k\}\qquad (k\in\mathbb{N}).
$$
对任意 $\delta\in(0,1)$，以及任意固定的 $k\in\mathbb{N}$，定义两侧证书点
$$
n_-(k,\delta):=\max\Bigl\{n\ge 1:\ k-\mu_e n>0\ \text{且}\ \frac{\mathrm{Var}(T_n)}{(k-\mu_e n)^2}\le\delta\Bigr\},
$$
$$
n_+(k,\delta):=\min\Bigl\{n\ge 1:\ \mu_e n-k>0\ \text{且}\ \frac{\mathrm{Var}(T_n)}{(\mu_e n-k)^2}\le\delta\Bigr\}.
$$
则由 Chebyshev 不等式与恒等式 $\mathbb{P}(\tau(k)>n)=\mathbb{P}(T_n<k)$，有
$$
\mathbb{P}\bigl(\tau(k)\le n_-(k,\delta)\bigr)\le\delta,
\qquad
\mathbb{P}\bigl(\tau(k)>n_+(k,\delta)\bigr)\le\delta,
$$
从而由并合界得到
$$
\mathbb{P}\bigl(n_-(k,\delta)<\tau(k)\le n_+(k,\delta)\bigr)\ge 1-2\delta.
$$
表 \ref{tab:real-input-40-hitting-time-deviation-certificate} 给出若干 $k$ 的可复现数值证书。
\end{remark}

% AUTO-GENERATED by scripts/exp_real_input_40_event_clock_deviation_certificate.py
\begin{table}[H]
\centering
\scriptsize
\setlength{\tabcolsep}{6pt}
\renewcommand{\arraystretch}{1.15}
\caption{事件时钟的击中时间 $\tau(k)$ 偏差证书（真实输入 40 状态核；$t=0$ Parry 平衡态）。设 $T_n=\sum_{t=0}^{n-1} e_t$ 且 $\mu_e=0.367376207875$，并定义 $\tau(k):=\inf\{n:\,T_n\ge k\}$。对给定 $\delta=0.05$，表中 $n_-(k,\delta)$ 为满足 $\mathbb{P}(\tau(k)\le n_-)\le\delta$ 的最大 $n$（左尾），$n_+(k,\delta)$ 为满足 $\mathbb{P}(\tau(k)>n_+)\le\delta$ 的最小 $n$（右尾），二者均由 $\mathbb{P}(|T_n-\mu_e n|\ge a)\le \mathrm{Var}(T_n)/a^2$（Chebyshev）与恒等式 $\mathbb{P}(\tau(k)>n)=\mathbb{P}(T_n<k)$ 推出。因此由并合界可得 $\mathbb{P}(n_-<\tau(k)\le n_+)\ge 1-2\delta$。}
\label{tab:real-input-40-hitting-time-deviation-certificate}
\begin{tabular}{r r r r r r}
\toprule
$k$ & $k/\mu_e$ & $n_-(k,\delta)$ & $\mathbb{P}(\tau(k)\le n_-)$ 上界 & $n_+(k,\delta)$ & $\mathbb{P}(\tau(k)>n_+)$ 上界\\
\midrule
5 & 13.610026705105 & 5.000000000000 & 0.045905808565 & 30.000000000000 & 0.046265197393\\
10 & 27.220053410210 & 15.000000000000 & 0.046951070365 & 47.000000000000 & 0.047457682717\\
20 & 54.440106820420 & 37.000000000000 & 0.049172745149 & 79.000000000000 & 0.049941439951\\
50 & 136.100267051051 & 107.000000000000 & 0.047513908181 & 172.000000000000 & 0.049433961235\\
\bottomrule
\end{tabular}
\end{table}


\begin{remark}[击中时间 $\tau(k)$ 的精确分位数区间（无近似）]
\label{prop:chi-event-clock-hitting-time-exact-quantiles}
在命题 \ref{prop:chi-event-clock-hitting-time-certificate} 的设定下，固定 $N\ge 1$。
将 $T_n$ 的分布视为有限状态马尔可夫链上有界边可观测量的加性泛函，则对每个 $n\le N$，$T_n$ 的精确分布可由状态空间上的递推
$$
\mathbb{P}(X_{t+1}=j,\,T_{t+1}=m) = \sum_i \mathbb{P}(X_t=i,\,T_t=m-e(i\to j))\,P_{ij}
$$
（并约定 $m-e(i\to j)<0$ 时该项为 $0$）在有限步内完全计算。
因此对任意 $k\in\{1,\dots,N\}$，击中时间 $\tau(k)=\inf\{n\ge 1:\,T_n\ge k\}$ 的分布函数满足恒等式
$$
\mathbb{P}(\tau(k)\le n)=\mathbb{P}(T_n\ge k),
$$
从而可在离散 $n=1,\dots,N$ 上给出任意分位数的\emph{精确}整数解。
表 \ref{tab:real-input-40-hitting-time-exact-quantiles} 报告若干 $k$ 的 $[q_{\mathrm{lo}},q_{\mathrm{hi}}]$ 分位数区间与参照点 $k/\mu_e$ 的最大偏差，作为“事件时间严格换时误差”在本文口径下的可审计证书。
\end{remark}

% AUTO-GENERATED by scripts/exp_real_input_40_hitting_time_exact_quantiles.py
\begin{table}[H]
\centering
\scriptsize
\setlength{\tabcolsep}{6pt}
\renewcommand{\arraystretch}{1.15}
\caption{事件时钟击中时间 $\tau(k)$ 的精确分位数区间（真实输入 40 状态核；$t=0$ Parry 平衡态）。设 $T_n=\sum_{t=0}^{n-1} e_t$，$\mu_e=0.367376207875$，$\tau(k)=\inf\{n\ge1:\,T_n\ge k\}$。表中 $n_{\mathrm{lo}}$ 与 $n_{\mathrm{hi}}$ 分别为 $\mathbb{P}(\tau(k)\le n_{\mathrm{lo}})\ge 0.05$ 与 $\mathbb{P}(\tau(k)\le n_{\mathrm{hi}})\ge 0.95$ 的最小整数解（精确分布递推，无近似）。其中 $k/\mu_e$ 为均值换时的参照点。}
\label{tab:real-input-40-hitting-time-exact-quantiles}
\begin{tabular}{r r r r r r}
\toprule
$k$ & $k/\mu_e$ & $n_{\mathrm{lo}}$ & $n_{\mathrm{med}}$ & $n_{\mathrm{hi}}$ & $\max\{|n_{\mathrm{lo}}-k/\mu_e|,|n_{\mathrm{hi}}-k/\mu_e|\}$\\
\midrule
5 & 13.610026705105 & 10 & 13 & 17 & 3.610026705105\\
8 & 21.776042728168 & 17 & 21 & 26 & 4.776042728168\\
10 & 27.220053410210 & 22 & 26 & 32 & 5.220053410210\\
20 & 54.440106820420 & 47 & 53 & 62 & 7.559893179580\\
50 & 136.100267051051 & 124 & 135 & 148 & 12.100267051051\\
\bottomrule
\end{tabular}
\end{table}


\begin{remark}[事件时间指数包络的误差预算：从 $k/\mu_e$ 到 $\tau(k)$]
\label{cor:chi-event-time-exponential-envelope-error-budget}
在引理 \ref{lem:chi-event-time-rescaling} 的设定下，取 $t=0$ 且令 $\rho:=\rho_{\mathrm{corr}}(0)=\varphi^{-1}$。
对任意 $k\in\mathbb{N}$，令 $\tau(k)=\inf\{n\ge 1:\,T_n\ge k\}$ 为事件时间的严格换时。
对给定分位数 $0<q_{\mathrm{lo}}<q_{\mathrm{hi}}<1$，令 $n_{\mathrm{lo}}(k),n_{\mathrm{hi}}(k)$ 分别为满足
$$
\mathbb{P}(\tau(k)\le n_{\mathrm{lo}}(k))\ge q_{\mathrm{lo}},
\qquad
\mathbb{P}(\tau(k)\le n_{\mathrm{hi}}(k))\ge q_{\mathrm{hi}}
$$
的最小整数解。
则由 $\rho\in(0,1)$ 的单调性，得到乘性误差预算：在概率至少 $q_{\mathrm{hi}}-q_{\mathrm{lo}}$ 的事件上，
$$
\rho^{\tau(k)}\in\Bigl[\rho^{n_{\mathrm{hi}}(k)},\,\rho^{n_{\mathrm{lo}}(k)}\Bigr]
=\rho^{k/\mu_e}\cdot\Bigl[\rho^{n_{\mathrm{hi}}(k)-k/\mu_e},\,\rho^{n_{\mathrm{lo}}(k)-k/\mu_e}\Bigr].
$$
当把事件阈值取为扭曲混合尺度 $k_{\mathrm{mix}}:=\lceil\mu_e\,\tau_{\mathrm{mix}}\rceil$ 时，上式将“角色扭曲的读出混合瓶颈”与“事件换时误差”在同一指数包络下对齐。
表 \ref{tab:real-input-40-event-time-rescaling-error-budget} 给出 $k_{\mathrm{mix}}$ 及若干参照 $k$ 的精确分位数与对应的乘性误差区间（可复现）。
\end{remark}

% AUTO-GENERATED by scripts/exp_real_input_40_event_time_rescaling_error_budget.py
\begin{table}[H]
\centering
\scriptsize
\setlength{\tabcolsep}{6pt}
\renewcommand{\arraystretch}{1.15}
\caption{事件时间换时误差预算：$\tau(k)$ 的精确分位数与指数包络扰动（真实输入 40 状态核；$t=0$ Parry 平衡态）。设 $\mu_e=0.367376207875$，$\rho=\rho_{\mathrm{corr}}(0)=0.618033988750$，$\tau(k)=\inf\{n\ge1:\,T_n\ge k\}$。表中 $n_{\mathrm{lo}},n_{\mathrm{hi}}$ 满足 $\mathbb{P}(\tau(k)\le n_{\mathrm{lo}})\ge 0.05$ 与 $\mathbb{P}(\tau(k)\le n_{\mathrm{hi}})\ge 0.95$。并给出 $\rho^{\tau(k)}$ 相对 $\rho^{k/\mu_e}$ 的乘性不确定性区间 $[\rho^{n_{\mathrm{hi}}-k/\mu_e},\,\rho^{n_{\mathrm{lo}}-k/\mu_e}]$（由 $\rho\in(0,1)$ 的单调性确定）。}
\label{tab:real-input-40-event-time-rescaling-error-budget}
\begin{tabular}{r r r r r r}
\toprule
$k$ & $k/\mu_e$ & $n_{\mathrm{lo}}$ & $n_{\mathrm{hi}}$ & $\max|n-k/\mu_e|$ & $\rho^{n_{\mathrm{hi}}-k/\mu_e} \sim \rho^{n_{\mathrm{lo}}-k/\mu_e}$\\
\midrule
5 & 13.610026705105 & 10 & 17 & 3.610026705105 & $0.195675963340\,\sim\,5.681342379773$\\
8 & 21.776042728168 & 17 & 26 & 4.776042728168 & $0.130992109849\,\sim\,9.957123630648$\\
10 & 27.220053410210 & 22 & 32 & 5.220053410210 & $0.100242119721\,\sim\,12.328965695215$\\
20 & 54.440106820420 & 47 & 62 & 7.559893179580 & $0.026307268894\,\sim\,35.883134057613$\\
50 & 136.100267051051 & 124 & 148 & 12.100267051051 & $0.003259139061\,\sim\,337.914056129726$\\
73 & 198.706389894535 & 185 & 213 & 14.293610105465 & $0.001029940541\,\sim\,731.924155675170$\\
\bottomrule
\end{tabular}
\end{table}


\begin{remark}[扭曲混合尺度处的量级对齐：事件换时误差为常数因子]
\label{cor:chi-event-time-envelope-vs-tau-mix}
在推论 \ref{cor:chi-event-time-exponential-envelope-error-budget} 的设定下，取 $k_{\mathrm{mix}}:=\lceil \mu_e\,\tau_{\mathrm{mix}}\rceil$。
令 $n_{\mathrm{lo}}(k),n_{\mathrm{hi}}(k)$ 为 $\tau(k)$ 的 $[q_{\mathrm{lo}},q_{\mathrm{hi}}]$ 精确分位数（命题 \ref{prop:chi-event-clock-hitting-time-exact-quantiles}）。
则在概率至少 $q_{\mathrm{hi}}-q_{\mathrm{lo}}$ 的事件上，成立
$$
\frac{\rho^{\tau(k_{\mathrm{mix}})}}{\rho^{\tau_{\mathrm{mix}}}}
\in\Bigl[\rho^{n_{\mathrm{hi}}(k_{\mathrm{mix}})-\tau_{\mathrm{mix}}},\,\rho^{n_{\mathrm{lo}}(k_{\mathrm{mix}})-\tau_{\mathrm{mix}}}\Bigr].
$$
表 \ref{tab:real-input-40-event-time-envelope-vs-tau-mix} 在 $k_{\mathrm{mix}}$ 的邻域给出上述比值区间的可复现数值读出，显示“事件时钟严格换时”对 $\rho^n$ 型指数包络的扰动在该对齐尺度上仅为常数倍。
因此，角色扭曲混合瓶颈 $\tau_{\mathrm{mix}}$ 与内生相关尺度之间的数量级分离（推论 \ref{cor:chi-event-clock-amplifies-separation}）不能归因于事件换时的不确定性。
\end{remark}

% AUTO-GENERATED by scripts/exp_real_input_40_event_time_rescaling_error_budget.py
\begin{table}[H]
\centering
\scriptsize
\setlength{\tabcolsep}{6pt}
\renewcommand{\arraystretch}{1.15}
\caption{事件时间指数包络在扭曲混合尺度处的对齐比较（真实输入 40 状态核；$t=0$ Parry 平衡态）。设 $\mu_e=0.367376207875$，$\rho=\rho_{\mathrm{corr}}(0)=0.618033988750$，$\tau_{\mathrm{mix}}=19.747340579192$，并记 $\rho^{\tau_{\mathrm{mix}}}=0.000074653451$。对 $k$ 的严格换时 $\tau(k)=\inf\{n\ge1:\,T_n\ge k\}$，令 $n_{\mathrm{lo}},n_{\mathrm{med}},n_{\mathrm{hi}}$ 分别为 $q_{\mathrm{lo}}=0.05$、$0.5$、$q_{\mathrm{hi}}=0.95$ 的精确分位数（由 $\mathbb{P}(\tau(k)\le n)=\mathbb{P}(T_n\ge k)$ 的有限步递推给出）。表中最后一列给出 $\rho^{\tau(k)}$ 在 $[q_{\mathrm{lo}},q_{\mathrm{hi}}]$ 区间上的乘性范围，相对于 $\rho^{\tau_{\mathrm{mix}}}$ 的比值区间 $[\rho^{n_{\mathrm{hi}}-\tau_{\mathrm{mix}}},\,\rho^{n_{\mathrm{lo}}-\tau_{\mathrm{mix}}}]$。}
\label{tab:real-input-40-event-time-envelope-vs-tau-mix}
\begin{tabular}{r r r r r r}
\toprule
$k$ & $k/\mu_e$ & $n_{\mathrm{lo}}$ & $n_{\mathrm{med}}$ & $n_{\mathrm{hi}}$ & $\rho^{n_{\mathrm{hi}}-\tau_{\mathrm{mix}}} \sim \rho^{n_{\mathrm{lo}}-\tau_{\mathrm{mix}}}$\\
\midrule
6 & 16.332032046126 & 12 & 15 & 20 & $0.885517817107\,\sim\,41.600488062556$\\
7 & 19.054037387147 & 14 & 18 & 23 & $0.209042400124\,\sim\,15.889972491312$\\
8 & 21.776042728168 & 17 & 21 & 26 & $0.049348216609\,\sim\,3.751113668551$\\
9 & 24.498048069189 & 19 & 24 & 29 & $0.011649533688\,\sim\,1.432797925722$\\
10 & 27.220053410210 & 22 & 26 & 32 & $0.002750081857\,\sim\,0.338237708491$\\
\bottomrule
\end{tabular}
\end{table}


\begin{remark}[分位数调整的事件相关尺度：$\tau_{\mathrm{eff}}^{(e)}(k)$]
\label{prop:chi-event-time-tau-eff}
在引理 \ref{lem:chi-event-time-rescaling} 的设定下，设步数时间相关满足指数包络 $\rho^n$，并记
$$
\tau_{\mathrm{corr}}(0):=\frac{1}{-\log\rho}.
$$
对事件阈值 $k\in\mathbb{N}$ 的击中时间 $\tau(k)$，定义一个与样本路径相关的“事件时间有效相关尺度”
$$
\tau_{\mathrm{eff}}^{(e)}(k):=\frac{k}{\tau(k)}\,\tau_{\mathrm{corr}}(0).
$$
若 $n_{\mathrm{lo}}(k),n_{\mathrm{med}}(k),n_{\mathrm{hi}}(k)$ 分别为 $\tau(k)$ 在 $q_{\mathrm{lo}},0.5,q_{\mathrm{hi}}$ 处的精确分位数，则由 $\tau(k)$ 的单调性立即得到分位数区间
$$
\tau_{\mathrm{eff}}^{(e)}(k)\in\Bigl[\frac{k}{n_{\mathrm{hi}}(k)}\tau_{\mathrm{corr}}(0),\,\frac{k}{n_{\mathrm{lo}}(k)}\tau_{\mathrm{corr}}(0)\Bigr]
$$
在概率至少 $q_{\mathrm{hi}}-q_{\mathrm{lo}}$ 的事件上成立。
\end{remark}

\begin{remark}[对齐到 $k_{\mathrm{mix}}$ 后的稳健分离]
\label{cor:chi-event-time-tau-eff-vs-tau-mix}
在命题 \ref{prop:chi-event-time-tau-eff} 的设定下取 $k_{\mathrm{mix}}:=\lceil\mu_e\tau_{\mathrm{mix}}\rceil$。
表 \ref{tab:real-input-40-event-time-tau-eff-vs-tau-mix} 给出 $k_{\mathrm{mix}}$ 邻域的 $\tau_{\mathrm{eff}}^{(e)}(k)$ 分位数区间，并与 $\tau_{\mathrm{mix}}$ 做比值对照。
该对照在数值上给出一条更强的审计结论：即使允许用 $\tau(k)$ 的分位数对事件时间进行非渐近、概率校准的换时，得到的事件时间相关尺度仍保持在常数阶，无法解释 $\tau_{\mathrm{mix}}$ 所体现的数量级混合瓶颈。
\end{remark}

% AUTO-GENERATED by scripts/exp_real_input_40_event_time_rescaling_error_budget.py
\begin{table}[H]
\centering
\scriptsize
\setlength{\tabcolsep}{6pt}
\renewcommand{\arraystretch}{1.15}
\caption{分位数调整后的事件时间相关尺度：$\tau_{\mathrm{eff}}^{(e)}(k)$ 与扭曲混合时间 $\tau_{\mathrm{mix}}$ 的对比（真实输入 40 状态核；$t=0$ Parry 平衡态）。设步数时间指数包络为 $\rho^n$，其中 $\rho=\rho_{\mathrm{corr}}(0)=0.618033988750$，$\tau_{\mathrm{corr}}(0)=1/(-\log\rho)=2.078086921235$。对事件阈值 $k$ 的击中时间 $\tau(k)$，以其精确分位数 $n_{\mathrm{lo}},n_{\mathrm{med}},n_{\mathrm{hi}}$ 诱导事件时间的有效相关尺度 $\tau_{\mathrm{eff}}^{(e)}(k):=(k/\tau(k))\,\tau_{\mathrm{corr}}(0)$。表中给出 $[q_{\mathrm{lo}},q_{\mathrm{hi}}]$ 区间内的尺度区间 $[\tau_{\mathrm{eff,lo}}^{(e)}(k),\tau_{\mathrm{eff,hi}}^{(e)}(k)]=[(k/n_{\mathrm{hi}})\tau_{\mathrm{corr}}(0),(k/n_{\mathrm{lo}})\tau_{\mathrm{corr}}(0)]$，其中 $q_{\mathrm{lo}}=0.05$，$q_{\mathrm{hi}}=0.95$，并比较 $\tau_{\mathrm{mix}}=19.747340579192$ 与该区间的比值。}
\label{tab:real-input-40-event-time-tau-eff-vs-tau-mix}
\begin{tabular}{r r r r r r}
\toprule
$k$ & $k/\mu_e$ & $\tau_{\mathrm{eff,lo}}^{(e)}(k)$ & $\tau_{\mathrm{eff,med}}^{(e)}(k)$ & $\tau_{\mathrm{eff,hi}}^{(e)}(k)$ & $\tau_{\mathrm{mix}}/\tau_{\mathrm{eff,hi}}^{(e)} \sim \tau_{\mathrm{mix}}/\tau_{\mathrm{eff,lo}}^{(e)}$\\
\midrule
6 & 16.332032046126 & 0.623426076371 & 0.831234768494 & 1.039043460618 & $19.005307600373\,\sim\,31.675512667289$\\
7 & 19.054037387147 & 0.632461236898 & 0.808144913814 & 1.039043460618 & $19.005307600373\,\sim\,31.223005343470$\\
8 & 21.776042728168 & 0.639411360380 & 0.791652160470 & 0.977923257052 & $20.193139325397\,\sim\,30.883624850607$\\
9 & 24.498048069189 & 0.644923527280 & 0.779282595463 & 0.984356962690 & $20.061158022616\,\sim\,30.619662245046$\\
10 & 27.220053410210 & 0.649402162886 & 0.799264200475 & 0.944584964198 & $20.905838360411\,\sim\,30.408492160597$\\
\bottomrule
\end{tabular}
\end{table}


\begin{remark}[期望级证书：$\mathbb{E}(\rho^{\tau(k)})$ 的严格区间]
\label{prop:chi-event-time-rho-tau-expectation}
在推论 \ref{cor:chi-event-time-exponential-envelope-error-budget} 的设定下，固定 $k\in\mathbb{N}$。
令 $F_n(k):=\mathbb{P}(\tau(k)\le n)=\mathbb{P}(T_n\ge k)$（命题 \ref{prop:chi-event-clock-hitting-time-exact-quantiles}），并定义离散概率质量
$$
p_n(k):=\mathbb{P}(\tau(k)=n)=F_n(k)-F_{n-1}(k)\qquad (n\ge 1),
$$
其中约定 $F_0(k)=0$。
则对任意截断 $n_{\max}\ge 1$，有
$$
\sum_{n=1}^{n_{\max}}\rho^n p_n(k)\ \le\ \mathbb{E}(\rho^{\tau(k)})\ \le\ \sum_{n=1}^{n_{\max}}\rho^n p_n(k)\ +\ \rho^{n_{\max}+1}\,\mathbb{P}(\tau(k)>n_{\max}).
$$
该区间完全由有限步递推得到的 $\{F_n(k)\}_{n\le n_{\max}}$ 决定，因而在本文口径下可作为无需外部不等式的“期望级误差预算”。
\end{remark}

\begin{remark}[对齐到 $k_{\mathrm{mix}}$ 的期望级对照]
\label{cor:chi-event-time-rho-tau-expectation-vs-tau-mix}
在命题 \ref{prop:chi-event-time-rho-tau-expectation} 的设定下取 $k_{\mathrm{mix}}:=\lceil\mu_e\tau_{\mathrm{mix}}\rceil$。
表 \ref{tab:real-input-40-event-time-rho-tau-expectation-vs-tau-mix} 给出 $k_{\mathrm{mix}}$ 邻域的 $\mathbb{E}(\rho^{\tau(k)})$ 严格区间，并与 $\rho^{\tau_{\mathrm{mix}}}$ 比较。
该对照在期望层面进一步表明：即使把事件换时的全部分布信息通过 $\mathbb{E}(\rho^{\tau(k)})$ 聚合，$k_{\mathrm{mix}}$ 邻域内的量级仍与 $\rho^{\tau_{\mathrm{mix}}}$ 处于可审计的分离区间，从而角色扭曲混合瓶颈并非换时噪声的产物。
\end{remark}

% AUTO-GENERATED by scripts/exp_real_input_40_event_time_rescaling_error_budget.py
\begin{table}[H]
\centering
\scriptsize
\setlength{\tabcolsep}{6pt}
\renewcommand{\arraystretch}{1.15}
\caption{事件换时的期望级证书：$\mathbb{E}(\rho^{\tau(k)})$ 的严格区间与 $\rho^{\tau_{\mathrm{mix}}}$ 对比（真实输入 40 状态核；$t=0$ Parry 平衡态）。设 $\mu_e=0.367376207875$，$\rho=\rho_{\mathrm{corr}}(0)=0.618033988750$，$\tau_{\mathrm{mix}}=19.747340579192$，并记 $\rho^{\tau_{\mathrm{mix}}}=0.000074653451$。对每个 $k$，用有限步递推得到 $F_n(k)=\mathbb{P}(\tau(k)\le n)=\mathbb{P}(T_n\ge k)$（$n\le 260$），从而 $p_n(k)=F_n(k)-F_{n-1}(k)$。由此计算 $\sum_{n=1}^{n_{\max}}\rho^n p_n(k)$，并用尾界 $\sum_{n>n_{\max}}\rho^n p_n(k)\le \rho^{n_{\max}+1}\,\mathbb{P}(\tau(k)>n_{\max})$ 给出 $\mathbb{E}(\rho^{\tau(k)})$ 的严格区间。}
\label{tab:real-input-40-event-time-rho-tau-expectation-vs-tau-mix}
\begin{tabular}{r r r r r r}
\toprule
$k$ & $k/\mu_e$ & $\mathbb{E}(\rho^{\tau(k)})_{\mathrm{lo}}$ & $\mathbb{E}(\rho^{\tau(k)})_{\mathrm{hi}}$ & $\mathbb{E}(\rho^{\tau(k)})/\rho^{\tau_{\mathrm{mix}}}$ & $\mathbb{P}(\tau(k)>n_{\max})$\\
\midrule
6 & 16.332032046126 & 0.001000488561 & 0.001000488561 & $13.401772341827\,\sim\,13.401772341827$ & 0.000000000000\\
7 & 19.054037387147 & 0.000294608145 & 0.000294608145 & $3.946343258878\,\sim\,3.946343258878$ & 0.000000000000\\
8 & 21.776042728168 & 0.000086750658 & 0.000086750658 & $1.162044840841\,\sim\,1.162044840841$ & 0.000000000000\\
9 & 24.498048069189 & 0.000025544602 & 0.000025544602 & $0.342175763032\,\sim\,0.342175763032$ & 0.000000000000\\
10 & 27.220053410210 & 0.000007521854 & 0.000007521854 & $0.100756946218\,\sim\,0.100756946218$ & 0.000000000000\\
\bottomrule
\end{tabular}
\end{table}


\begin{remark}[期望有效指数与偏移：$\tau^{(E)}(k)$ 与 $\Delta^{(E)}(k)$]
\label{prop:chi-event-time-tauE-delta}
在命题 \ref{prop:chi-event-time-rho-tau-expectation} 的设定下，假设 $0<\rho<1$。
定义事件换时的期望有效指数
$$
\tau^{(E)}(k):=\frac{\log\mathbb{E}(\rho^{\tau(k)})}{\log\rho},
$$
以及其相对线性重标定 $k/\mu_e$ 的偏移
$$
\Delta^{(E)}(k):=\tau^{(E)}(k)-\frac{k}{\mu_e}.
$$
若已知 $\mathbb{E}(\rho^{\tau(k)})\in[E_{\mathrm{lo}},E_{\mathrm{hi}}]$ 为严格区间，则由 $\log$ 的单调性与 $\log\rho<0$，得到
$$
\tau^{(E)}(k)\in\Bigl[\frac{\log E_{\mathrm{hi}}}{\log\rho},\,\frac{\log E_{\mathrm{lo}}}{\log\rho}\Bigr],
\qquad
\Delta^{(E)}(k)\in\Bigl[\frac{\log E_{\mathrm{hi}}}{\log\rho}-\frac{k}{\mu_e},\,\frac{\log E_{\mathrm{lo}}}{\log\rho}-\frac{k}{\mu_e}\Bigr].
$$
因此 $\Delta^{(E)}(k)$ 可被视为“期望层面”的事件换时偏移证书：它把 $\tau(k)$ 的全分布信息通过 $\mathbb{E}(\rho^{\tau(k)})$ 聚合为一个可审计的有效指数偏移。
\end{remark}

\begin{remark}[对齐到 $k_{\mathrm{mix}}$ 的期望偏移证书]
\label{cor:chi-event-time-tauE-delta-vs-tau-mix}
在命题 \ref{prop:chi-event-time-tauE-delta} 的设定下取 $k_{\mathrm{mix}}:=\lceil\mu_e\tau_{\mathrm{mix}}\rceil$。
表 \ref{tab:real-input-40-event-time-tauE-delta-vs-tau-mix} 给出 $k_{\mathrm{mix}}$ 邻域内的 $\tau^{(E)}(k)$ 与 $\Delta^{(E)}(k)$ 的严格区间。
该对照把“事件时间 $k$ 与步数时间指数包络”之间的换时误差在期望层面压缩为可复现的偏移量，并显示该偏移在对齐尺度处仅为常数阶。
\end{remark}

% AUTO-GENERATED by scripts/exp_real_input_40_event_time_rescaling_error_budget.py
\begin{table}[H]
\centering
\scriptsize
\setlength{\tabcolsep}{6pt}
\renewcommand{\arraystretch}{1.15}
\caption{事件换时的期望有效指数与偏移：定义 $\tau^{(E)}(k):=\log\mathbb{E}(\rho^{\tau(k)})/\log\rho$，并记 $\Delta^{(E)}(k):=\tau^{(E)}(k)-k/\mu_e$（真实输入 40 状态核；$t=0$ Parry 平衡态）。设 $\mu_e=0.367376207875$，$\rho=\rho_{\mathrm{corr}}(0)=0.618033988750$，$\tau_{\mathrm{mix}}=19.747340579192$，并用 $n_{\max}=260$ 的有限递推与尾界得到 $\mathbb{E}(\rho^{\tau(k)})$ 的严格区间（见表 \ref{tab:real-input-40-event-time-rho-tau-expectation-vs-tau-mix}）。由 $\log$ 单调性与 $\log\rho<0$，可把该区间推出 $\tau^{(E)}(k)$ 与 $\Delta^{(E)}(k)$ 的严格区间。}
\label{tab:real-input-40-event-time-tauE-delta-vs-tau-mix}
\begin{tabular}{r r r r r r}
\toprule
$k$ & $k/\mu_e$ & $\tau^{(E)}(k)_{\mathrm{lo}}$ & $\tau^{(E)}(k)_{\mathrm{hi}}$ & $\Delta^{(E)}(k)_{\mathrm{lo}}$ & $\Delta^{(E)}(k)_{\mathrm{hi}}$\\
\midrule
6 & 16.332032046126 & 14.353900876823 & 14.353900876823 & -1.978131169304 & -1.978131169304\\
7 & 19.054037387147 & 16.894564896219 & 16.894564896219 & -2.159472490928 & -2.159472490928\\
8 & 21.776042728168 & 19.435250893884 & 19.435250893884 & -2.340791834284 & -2.340791834284\\
9 & 24.498048069189 & 21.975944888338 & 21.975944888338 & -2.522103180851 & -2.522103180851\\
10 & 27.220053410210 & 24.516641780697 & 24.516641780697 & -2.703411629513 & -2.703411629513\\
\bottomrule
\end{tabular}
\end{table}


\paragraph{指数包络的离散审计证书}
在 $t=0$ 的 Parry 平衡态下，脚本输出的自相关序列 $\{\mathrm{Corr}(n)\}_{n=0}^N$ 与谱隙包络 $r^n$（$r=\rho_{\mathrm{corr}}(0)$）允许构造一个完全可复现的离散证书：
\[
C_N:=\max_{0\le n\le N}\frac{|\mathrm{Corr}(n)|}{r^n}.
\]
